\section{武周精英、佛教与书写政治} \label{sec:wu-zhou-elites-buddhism-writing}  \noindent	extbf{人物并不脱离网络。} 本节围绕立后、废立、建周、边地危机与复唐展开，账本入口为ST-655-WU-EMPRESS、ST-680-LI-XIAN、ST-690-ZHOU-FOUNDATION、ST-695-KHITAN-REBELLION、ST-705-TANG-RESTORATION。传记最容易把复杂关系收束为忠奸、功罪或个人性格；真正的行动却依靠亲属、门生、同僚、军府、寺院、使节与书吏的长期联结。一个人在本纪中获得显名，往往因为其行动同继承、战争或制度转折相接；这并不表示他可以代表整个群体。人物材料的价值，是让我们追问资源和信息如何在网络中流动，而不是把网络重新缩成少数“伟人”的意志。  \noindent	extbf{精英的可见性。} 墓志、家传、诗文、奏议和官修列传使仕宦家族与宗教名流特别可见。它们能够钉住姓名、籍贯、任官、迁居和丧葬的某些节点，也会通过追赠、套语和家族自我塑造重排记忆。研究者不应把志文中的德行直接当作政治事实，更不能用精英的教育、消费和宗教赞助替代城市居民、军户、农民和流动者的共同经验。材料的偏向必须成为叙事的一部分。  \noindent	extbf{佛教、道教与地方组织。} 寺院、道观、译场、造像与写经使宗教社群成为知识、慈善、财产和声望的组织者。国家的保护、诏令、限制或废并会改变这些组织的资源，却不能简单决定信仰如何在家庭、市场和乡村中实践。僧传追求典范，道教与佛教文本有各自的传承目的，题记又常只保存捐施者；将它们同契约、墓志和行政文书并读，才能说明宗教既可能服务国家礼仪，也可能提供跨地域的互助与记忆。  \noindent	extbf{文学与制度。} 诗文、笔记、表章和史论并非政治之外的闲暇。科举、学校、官职、宗教赞助和城市市场决定谁有时间、纸张和听众；易代和战乱又改变何种文体能够保存。文学文本的修辞不等于事件实录，但它能够显示作者如何理解道路、城市、战争、宗教和地方名望。把文学还原到生产、流通和制度条件中，可以避免既把它当透明史料，也把它隔离为无关社会的审美对象。  \noindent	extbf{外来者与跨境关系。} 使节、译者、商人、僧侣和军人沿草原、绿洲、海港与河道移动。中原史书的“朝贡”“归附”是政治语言，不能直接说明稳定的支配；外部史书、碑铭和多语文书有不同的纪年和目的，也不必然补足全部空白。较可靠的做法，是确认具体接触、战争、婚盟或翻译的时间，再说明材料无法证实的规模、动机和持续性。跨境交流既产生知识与物资的流动，也受战争、税收、身份分类和不平等约束。  \noindent	extbf{史料限度。} 两《唐书》、两《五代史》、《资治通鉴》、宗教传记、碑志、文集与考古材料提供互相不同的窗口。它们足以支持对事件年序和网络条件的分析，不能让我们精确统计信徒、读者、商人或地方追随者。叙事应把可确认的任命、文本和行动，与推测性的社会后果明确分开；这种节制不是减少解释，而是让人物、宗教、文学和外国关系都保留其历史复杂性。  